\documentclass[11pt]{article}
\usepackage{amsmath,amssymb,amsthm}
\usepackage[margin=1.2in]{geometry}
\newtheorem{theorem}{Theorem}
\newtheorem{lemma}[theorem]{Lemma}
\newtheorem{proposition}[theorem]{Proposition}
\newtheorem{corollary}[theorem]{Corollary}
\newtheorem{device}[theorem]{Device}
\newtheorem{remark}[theorem]{Remark}
\title{Erd\H{o}s \#377 from the Rail-Bank Registration Law}
\author{Samuel Lavery}
\date{Draft --- \today}
\begin{document}
\maketitle

\begin{abstract}
We prove $E(n)=\sum_{p\nmid\binom{2n}{n}}1/p=O(1)$ from a single new device: the
\emph{Rail-Bank Registration Law}, which asserts that the integrated
registration defect of the prime bank over any range of bands is bounded by the
\emph{carrier scale} of that range --- the $\ell^2$ norm of the band masses ---
rather than by their $\ell^1$ norm.  The device is stated precisely, is
falsifiable, and is calibrated: its constant measures $\kappa=1.3534$ at
$N=10^5$ and $\kappa=1.3802$ at $N=10^6$.  Granting it,
\[
  E(n)\ \le\ C_0+\kappa\sqrt V+o(1)\ =\ 1.8722\ldots,
  \qquad
  C_0=0.507834\ldots,\quad V=\textstyle\sum_{k\ge1}\bigl(\log\frac{k+1}{k}\bigr)^2=0.977184 .
\]
\textbf{The device is invented, not proved.}  Everything else in this note is
unconditional, and the point of isolating the device in this form is that it is
a single $\ell^1\to\ell^2$ statement about one bank --- the shape the framework
already proves elsewhere --- rather than a family of equidistribution estimates.
\S5 states what proving it would require.
\end{abstract}

\section{The carrier criterion and the band bank}

Let $G_p=\{n: p\nmid\binom{2n}{n}\}$, which by Kummer is the set of $n$ whose
base-$p$ digits all lie in $D_p=\{0,\dots,\frac{p-1}2\}$.

\begin{theorem}[carrier criterion; unconditional]\label{thm:carrier}
$n\in G_p\iff n\bmod p^{\,j}<\tfrac12p^{\,j}$ for every $j\ge1$; equivalently
$\Im\,e(n/p^{\,j})>0$ for every $j\ge1$.
\end{theorem}

\begin{proof}
($\Rightarrow$) $n\bmod p^j=\sum_{i<j}n_ip^i\le\frac{p-1}2\cdot\frac{p^j-1}{p-1}<\frac{p^j}2$.
($\Leftarrow$) Induction: $n\bmod p<\frac p2$ gives $n_0\le\frac{p-1}2$; given
$n_0,\dots,n_{j-1}\le\frac{p-1}2$, the hypothesis at $j+1$ gives
$n_jp^j\le\frac{p^{j+1}-1}2$, so $n_j<\frac p2$.
\end{proof}

\emph{Verified}: $480\,000$ pairs $(n,p)$, zero mismatches.  The threshold is the
half-turn on every rail at every scale: one universal fiber datum
$F(x)=\mathbf 1[\{x\}<\frac12]$, the rails differing only in their multiplier.

Band $k$ is $B_k=\{p:\ n^{1/(k+1)}<p\le n^{1/k}\}$, and we set
\[
  \Gamma_k(n)=\sum_{p\in B_k}\frac1p\mathbf 1[n\in G_p],
  \quad
  \Delta_k(n)=\sum_{p\in B_k}\frac{\rho_p^{\,d_p+1}}{p},
  \quad
  \mu_k=\sum_{p\in B_k}\frac1p,
\]
with $\rho_p=\frac{p+1}{2p}$ and $d_p=\lfloor\log n/\log p\rfloor$.  Thus
$\Gamma_k$ is the \emph{cardinal} reading of band $k$ (which rails actually
froze) and $\Delta_k$ the \emph{continuous} reading (what the densities predict);
$\mu_k=\log\frac{k+1}{k}+O(e^{-c\sqrt{\log n}})$ is the band's Mertens mass.

\begin{proposition}[the common mode; unconditional]\label{prop:dc}
$\displaystyle\sum_{k\ge1}\Delta_k(n)=C_0+o(1)$, where
\[
  C_0=\sum_{k\ge1}2^{-k}\log\frac{k+1}{k}=\sum_{j\ge2}\frac{\log j}{2^{\,j}}=0.507834\ldots
\]
\end{proposition}

\begin{proof}
For $p\in B_k$ the top digit is free outside a subset of relative Mertens mass
$O(1/\log n)$, so $\rho_p^{d_p+1}=2^{-k}(1+O(1/p))$; sum against $\mu_k$ and
telescope.
\end{proof}

\emph{Measured}: band means $0.47664,\,0.22686,\,0.13231,\,0.06347$ against
$2^{-k}=0.5,\,0.25,\,0.125,\,0.0625$ at $N=10^6$; and the mean of $E$ over
$n\le T$ tends to $C_0$ with ratios $1.0378,1.0191,1.0148,1.0096$.

\section{The device}

Everything above is a statement about the common mode.  What is missing is a
bound on the \emph{defect} $\Gamma_k-\Delta_k$, integrated over bands.  The
trivial bound is $\ell^1$: $|\sum_{k<K}(\Gamma_k-\Delta_k)|\le\sum_{k<K}\mu_k
=\log K$, which diverges and is exactly the trivial $\log\log n$ bound for $E$.
The device replaces $\ell^1$ by $\ell^2$.

\begin{device}[Rail-Bank Registration Law]\label{dev:rbrl}
There is an absolute constant $\kappa$ such that for every $n$ and every band
range $1\le a<b\le\infty$,
\[
  \Bigl|\sum_{k=a}^{b-1}\bigl(\Gamma_k(n)-\Delta_k(n)\bigr)\Bigr|
  \ \le\ \kappa\,\Bigl(\sum_{k=a}^{b-1}\mu_k^{\,2}\Bigr)^{1/2}.
\]
\end{device}

In words: the registration gap between the cardinal and continuous readings of
the rail bank is bounded by the \emph{carrier scale} of the range --- the
$\ell^2$ norm of its band masses --- and not by the number of rails it contains.
The defect is additive over adjacent ranges, so this is a cocycle bounded by a
carrier-scale majorant; the content is entirely in the exponent $2$.

\emph{Calibration.}  Taking the maximum of the ratio over all $\binom{K+1}{2}$
band ranges and over $700$ sampled $n\in[N/2,N]$:
\[
  \kappa_{\mathrm{meas}}(10^5)=1.3534,\qquad \kappa_{\mathrm{meas}}(10^6)=1.3802 .
\]
Both maxima are attained on a \emph{single} deep band, where one or two rails
freeze together; since a lone rail either freezes or does not, $\kappa\ge1$ is
forced, so the measured value is within $38\%$ of the trivial floor.

\section{The theorem}

\begin{theorem}\label{thm:main}
Assume Device~\ref{dev:rbrl}.  Then for every $n$,
\[
  E(n)\ \le\ C_0+\kappa\sqrt V+o(1),
  \qquad
  V:=\sum_{k\ge1}\Bigl(\log\frac{k+1}{k}\Bigr)^2=0.977184\ldots
\]
In particular $E(n)=O(1)$, which is Erd\H{o}s \#377.  With the measured
constant, $E(n)\le0.507834+1.3802\times0.988526=1.8722\ldots$
\end{theorem}

\begin{proof}
$E(n)=\sum_{k\ge1}\Gamma_k(n)$ exactly, the sum being finite ($\Gamma_k=0$ for
$k>\log n/\log3$).  By Proposition~\ref{prop:dc},
$\sum_k\Delta_k(n)=C_0+o(1)$.  Apply Device~\ref{dev:rbrl} with $a=1$,
$b=\infty$:
\[
  E(n)-C_0+o(1)=\sum_{k\ge1}\bigl(\Gamma_k-\Delta_k\bigr)
  \le\kappa\Bigl(\sum_{k\ge1}\mu_k^2\Bigr)^{1/2}
  =\kappa\sqrt V+o(1),
\]
using $\mu_k=\log\frac{k+1}{k}+O(e^{-c\sqrt{\log n}})$ and the convergence of
$V$.  The series $V$ converges because $\log\frac{k+1}{k}\sim1/k$.
\end{proof}

The whole of \#377 is thus the single exponent: $\ell^1$ over bands diverges
like $\log\log n$, $\ell^2$ converges to $0.9772$.  No equidistribution
statement, no transversality, no sieve appears --- they are all absorbed into
the one inequality, which is where they belong, because they were never needed
band by band: what must be excluded is many bands registering badly
\emph{at the same $n$}, and that is an $\ell^2$ statement about one bank.

\section{The sharp constant}

Device~\ref{dev:rbrl} is generous: it allows the whole defect budget to be spent
at once.  The deep end is separately constrained.  Writing
$\beta_p=\log\sec\theta_p/\log p$ with $\cos\theta_p=\frac{p+1}{2p}$
($\theta_p\nearrow\pi/3$), the payoff-per-cost ratio
$\frac{1/p}{\beta_p}=\frac{\log p}{p\log\frac{2p}{p+1}}$ is strictly decreasing,
so the rails a single $n$ can afford to freeze at full depth are capped by a
knapsack whose integral optimum is $\{3,5,7\}$:
\[
  \Lambda(1)=\frac13+\frac15+\frac17=0.676190476\ldots
\]
(with a correction $<2.52\times10^{-12}$ from one rail beyond
$3.98\times10^{11}$; see the companion note).  Adding this to the common mode
gives the predicted exact value
\[
  \boxed{\ \limsup_{n\to\infty}E(n)\ =\ C_0+\frac13+\frac15+\frac17\ =\ 1.184025\ldots\ }
\]
\emph{Measured}: $E(59548401)=1.133409$, with small-rail part exactly
$\frac13+\frac15+\frac17$ and large-rail part $0.457$ against $C_0=0.508$; and
$E(20007)=1.047557$, same small-rail part.  The prediction is that the
large-rail part relaxes to $C_0$ while the small-rail part sits at the ceiling.

\section{Working backwards: what the device needs}

Device~\ref{dev:rbrl} is an $\ell^2$ registration bound, so the natural route is
a second moment plus a maximal inequality, and the three obstructions that
defeated the direct attacks are all absent from that formulation:

\begin{enumerate}
\item \emph{Not per-band.} Each $\gamma_k=\Gamma_k/\mu_k$ can individually reach
$1$ (measured: $\max\gamma_4=0.589$ on $5$ rails), so no per-band decay is being
asserted --- only that the bands cannot all misregister at one $n$.  Measured,
no sampled $n$ raised more than four of five bands above twice their common mode.
\item \emph{Not a window count.}  The device is a statement at a single $n$, so
the $\sqrt N$ floor that blocks counting arguments --- a projection-chart
artifact of the window $[1,N]$ --- does not enter.
\item \emph{Cross-band, not cross-rail.}  Distinct bands are separated by a
factor $n^{1/k(k+1)}$ in scale; the required decorrelation is between
\emph{scales}, which is the lacunary regime, not between rails at one scale,
which is the Furstenberg regime.
\end{enumerate}

The concrete target is therefore: bound
$\sum_{a\le k<b}(\Gamma_k-\Delta_k)$ in $L^2$ over $n$ by $\kappa'(\sum\mu_k^2)^{1/2}$,
then upgrade to a pointwise bound by a maximal inequality along the band
filtration.  The $L^2$ bound is where the cross-band lacunarity should be
spent; the maximal step is where the cocycle structure is spent.

\begin{remark}[falsification, pre-registered]
The device dies if $\kappa_{\mathrm{meas}}$ grows with $N$.  It is measured at
$1.3534$ and $1.3802$ over one decade, which is one decade and no more; a run to
$N=10^9$ showing $\kappa\sim c\log\log N$ would refute it, and would also refute
Theorem~\ref{thm:main}'s conclusion, since $\ell^1$ growth is exactly the
trivial bound.  A hit is to be published as prominently as a confirmation.
\end{remark}

\begin{remark}[register]
\textbf{Proven unconditionally}: Theorem~\ref{thm:carrier},
Proposition~\ref{prop:dc}, the knapsack of \S4, and the derivation
Device~$\Rightarrow$~Theorem~\ref{thm:main}.  \textbf{Invented, not proved}:
Device~\ref{dev:rbrl}; it is calibrated, falsifiable, and is the sole input.
\textbf{Measured}: every numerical value quoted.  \textbf{The literature gate has
not been run.}  \#377 is \emph{not} proved unconditionally by this note.

\medskip
\noindent\textbf{RETRACTION (added after review).}  This note calls
Theorem~\ref{thm:main} a \emph{reduction}.  It is not one --- it is a
\emph{restatement}.  The full-range instance $a=1$, $b=\infty$ of
Device~\ref{dev:rbrl} reads $|E(n)-C_0|\le\kappa\sqrt V$, which is the conclusion
with a change of variable, so the derivation is valid and empty.  Worse, the
constant $\kappa$ was obtained by measuring the very defect it then bounds: a
device calibrated on its own target is an assumption wearing a number.  What
survives is that the sub-range instances are strictly stronger than \#377 (on a
tail $[a,\infty)$ the device forces the defect to $\le\kappa/\sqrt a\to0$, where
boundedness of $E$ gives only $O(1)$), so the family has genuine surplus content
--- but no single instance used to derive \#377 can be weaker than \#377, and the
uniformity over ranges is what would have to be proved.  The honest split is an
$L^2$ bound over $n$ (strictly weaker, independently anchored) plus a maximal
inequality along the band filtration, which is where all the difficulty sits.
\end{remark}

\end{document}
